\documentclass[10pt,a4paper]{amsart}
\usepackage[margin=19mm]{geometry}
\usepackage{amsmath,amssymb,mathtools}
\usepackage[scr=rsfs,cal=euler]{mathalfa}
\usepackage{xfrac}
\usepackage[unicode,hidelinks,bookmarks=false]{hyperref}
\newcommand{\ie}{i.e.\ }
\newcommand{\cf}{cf.\ }
\newcommand{\resp}{respectively\ }
\newcommand{\Subsection}{\subsection}
\providecommand{\nobreakdash}{\nobreak\hbox{-}}
\setlength{\emergencystretch}{2em}
\multlinegap0pt
\usepackage{amssymb}
\usepackage{mathtools}
\usepackage{xcolor}
\newtheorem{lemma}{Lemma}
\newtheorem{claim}{Claim}
\DeclareMathOperator*{\argmin}{\arg\!\min}
\newcommand{\cC}{\mathcal{C}}
\newcommand{\cG}{\mathcal{G}}
\newcommand{\norm}[1]{\left\| #1\right\|_{2}}
\newcommand{\paren}[1]{\left( #1\right)}
\title{A Short and Unified Convergence Analysis of the SAG, SAGA, and IAG Algorithms}
\author{Feng Zhu, Robert W. Heath Jr.,, Aritra Mitra}
\date{Source: arXiv:2602.05304v2 (CC BY 4.0)}
\begin{document}
\maketitle

\noindent\emph{Source and licence.} A Short and Unified Convergence Analysis of the SAG, SAGA, and IAG Algorithms. Version arXiv:2602.05304v2 (CC BY 4.0), distributed by arXiv under the Creative Commons Attribution 4.0 International licence, http://creativecommons.org/licenses/by/4.0/, as recorded in the arXiv licence metadata for this exact version and shown on the abstract page https://arxiv.org/abs/2602.05304v2. Full licence notice: \texttt{licenses/arxiv-2602.05304-CC-BY-4.0.txt}.\par\smallskip

\noindent\textit{Editorial scope.} Counted unit: the article's claim claim1 (source0.tex lines 747--753). The article's complete original proof of it is delivered with it (source0.tex lines 754--801, one \texttt{proof} environment of 48 lines). No mathematics is written, repaired or re-derived: the statement and the proof are the article's own lines, in the article's own notation, and no retained line is substituted or re-flowed.  Retained uncounted as the source-local context the proof needs: lines 442--448; lines 543--546.  Nothing was omitted from any retained span, so the package carries no omission record.  No retained line contains a citation, so the wrapper carries no bibliography and no bibliography entry.  No retained line contains a figure environment, an image-inclusion command or drawing code, so no image is delivered and the package declares no asset dependency.  Editorial formatting only, in addition: an AMS \texttt{amsart} preamble that restates the article's own declarations and macros used by the retained lines, namely the environments \texttt{lemma}, \texttt{claim} on separate counters, together with the standard packages those lines need.  The retained environments print in this document as Lemma 1, Claim 1 in order of appearance, whereas the article prints the same items with the numbering of its own full document.  The complete native source of the article as distributed in the arXiv e-print for this exact version is bundled unchanged as \texttt{sources/194025/source0.tex}.
\begin{lemma}[Gradient Error]\label{lem:e_k} On event $\cG$, the gradient error $e_k$ satisfies the following for 
both SAG and SAGA:
\begin{equation}
    \norm{e_k}^2\leq 4L^2\tau\alpha^2  U_k,\ \  \forall k\geq \tau,\ \  \textup{with } U_k:=\sum_{j=1}^\tau \norm{g_{k-j}}^2.
\nonumber
\end{equation}
\end{lemma}

\begin{lemma}[Burn-In Effects]\label{lem:burnin}
Define $B:=\max\{\norm{x^*}, \norm{x_1^*},\cdots\norm{x_N^*}\}$, where $x_i^* \in \argmin_{x \in \mathbb{R}^d} f_i(x)$. With $\alpha = 1/(16 L \tau)$, the following is true: 
    $V_\tau\leq 3LB^2.$
\end{lemma}

\begin{claim}[Bounded Iterates]\label{claim1}
Fix $0\leq k\leq \tau$, we claim that the following holds for both SAG and SAGA:
\begin{equation}
    \norm{x_k}\leq B,
\end{equation}
where $B$ is as defined in Lemma~\ref{lem:burnin}. 
\end{claim}

\begin{proof}
We prove this claim by induction and first focus on the analysis of SAG. The base case holds trivially by initializing $x_0=0$. Suppose that this claim holds up to iteration $0\leq k\leq \tau-1$. Then for iteration $k+1$, we have
\begin{equation}
\begin{aligned}
        x_{k+1}&\overset{(a)}{=}x_k-\frac{\alpha}{N}\paren{\sum_{i\neq i_k}\nabla f_i(x_{\tau_{i,k}})+\nabla f_{i_k}(x_k)}\\
        &\overset{(b)}{=}x_k-\alpha\nabla f(x_k) -\frac{\alpha}{N}\paren{\sum_{i\neq i_k}\nabla f_i(x_{\tau_{i,k}})+\nabla f_{i_k}(x_k)-\sum_{i=1}^N\nabla f_i(x_k)}\\
        &=x_k-\alpha\nabla f(x_k)-\frac{\alpha}{N}\sum_{i\neq i_k}\paren{\nabla f_i(x_{\tau_{i,k}}) - \nabla f_i(x_k)},
\end{aligned}
\end{equation}
where $(a)$ uses the definition of $g_k^\textbf{SAG}$, and $(b)$ holds by adding and subtracting $\nabla f(x_k)$.

Taking the 2-norm on both sides and using triangle inequality yields
\begin{equation}
\begin{aligned}
    \norm{x_{k+1}}&\leq \norm{x_k}+\alpha\norm{\nabla f(x_k)-\nabla f(x^*)}+\frac{\alpha}{N}\sum_{i\neq i_k}\norm{\nabla f_i(x_{\tau_{i,k}}) - \nabla f_i(x_k)}\\
    &\overset{(a)}{\leq} \norm{x_k}+\alpha L \norm{x_k-x^*}+\frac{\alpha}{N}\sum_{i\in\cC_k, i\neq i_k}L\norm{x_{\tau_{i,k}} - x_k}+\frac{\alpha}{N}\sum_{i\in[N]\backslash\cC_k, i\neq i_k}\norm{\nabla f_i(x_k)-\nabla f_i(x_i^*)}\\
    &\overset{(b)}{\leq} \norm{x_k}+\alpha L \paren{\norm{x_k} + \norm{x^*}}+\frac{\alpha}{N}\sum_{i\in\cC_k, i\neq i_k}L\paren{\norm{x_{\tau_{i,k}}} + \norm{x_k}}+\frac{\alpha}{N}\sum_{i\in [N]\backslash\cC_k, i\neq i_k}L\paren{\norm{x_{k}} + \norm{x_i^*}}\\
    &\overset{(c)}{\leq} (1+L\alpha)\norm{x_k}+3LB\alpha,\label{eqn:claim_1}
\end{aligned}
\end{equation}
where $(a)$ uses smoothness, and $\cC_k$ denotes the set of component indices that have been sampled at least once up to iteration $k$. The reason for this definition is that in the first inequality, $\nabla f_i(x_{\tau_{i,k}})$ might be set to 0 since it is possible that component $i$ has never been accessed up to iteration $k$. Inequality $(b)$ holds due to smoothness and the fact that $\nabla f_i(x_i^*)=0$. Inequality $(c)$ uses the definition of $B$ and the induction hypothesis up to iteration $k$. 

Iterating~\eqref{eqn:claim_1} for $k+1$ steps from $k=0$ yields
\begin{equation}
\begin{aligned}
    \norm{x_{k+1}}&\leq (1+L\alpha)^{k+1}\norm{x_0} + 3LB\alpha \sum_{j=0}^{k}(1+L\alpha)^{k-j}\\
    &\leq 3LB\alpha \cdot \tau (1+L\alpha)^\tau\leq 4LB\tau\alpha\leq \frac{B}{4}\leq B,
\end{aligned}
\end{equation}
where we use the fact that $x_0=0$, $k+1\leq \tau$, $1+x\leq e^x$, and $\alpha\leq 1/(16L\tau)$. The proof is then complete for the SAG case.

The proof for the case of SAGA carries through similarly as in the proof of Lemma~\ref{lem:e_k}. Specifically, the one-step descent of SAGA can be written as
\begin{equation}
\begin{aligned}
    x_{k+1} &= x_k - \alpha\nabla f(x_k) - \frac{\alpha}{N}\sum_{i=1}^N\paren{\nabla f_i(x_{\tau_{i,k}})-\nabla f_i(x_k)}-\alpha\paren{\nabla f_{i_k}(x_k)-\nabla f_{i_k}(x_{\tau_{{i_k},k}})}.
\end{aligned}
\end{equation}
Taking the 2-norm on both sides and using triangle inequality yields
\begin{equation}
\begin{aligned}
    \norm{x_{k+1}}&\leq \norm{x_k}+\frac{\alpha}{N}\sum_{i=1}^N\norm{\nabla f_i(x_{\tau_{i,k}}) - \nabla f_i(x_k)}+ \alpha\norm{\nabla f_{i_k}(x_k)-\nabla f_{i_k}(x_{\tau_{{i_k},k}})} +\alpha\norm{\nabla f(x_k)-\nabla f(x^*)}\\
    &\leq (1+L\alpha)\norm{x_k}+5LB\alpha.\label{eqn:claim_2}
\end{aligned}
\end{equation}
The rationale is similar to the SAG case and is hence omitted, where one needs to break the analysis into two cases: (i) component $i$ has been sampled before iteration $k$, and (ii) it has never been sampled before iteration $k$. As has been shown for the SAG analysis, both cases yield the exact same bound. Iterating~\eqref{eqn:claim_2} for $k+1$ steps from $k=0$ yields the same bound. The proof is then complete.


\end{proof}

\end{document}
